\documentclass[twoside,a4paper]{article}
\usepackage[utf8]{inputenc}
\usepackage[a4paper,bindingoffset=0.5cm,inner=2.5cm,outer=2cm,top=2.5cm,bottom=2.5cm]{geometry}
%\usepackage[a4paper,bindingoffset=1cm,inner=2cm,outer=2cm,top=2.5cm,bottom=2.5cm]{geometry} %,showframe
\usepackage{helvet}
\usepackage[T1]{fontenc}
\renewcommand{\familydefault}{\sfdefault}
% \usepackage[german,ngerman]{babel}
\usepackage[english]{babel}
\usepackage{graphicx}
\usepackage{amsmath,amsthm,amstext,amssymb,bm}
\usepackage{mathtools}
\usepackage{xcolor,color}
\usepackage{pifont}
\usepackage{array}
\usepackage{upgreek}
\usepackage{enumerate}
\usepackage{pgf}
\usepackage{mathrsfs}
\usepackage{subfig}
\usepackage{tabularx}
\usepackage[numbers]{natbib}
\usepackage{multirow}
\usepackage{makecell}
%%% Title page
\usepackage{common/titlePageST}
%%% Appendix in TOC
\usepackage[toc,page]{appendix}
%%% Setstrech on title page
\usepackage{setspace}
%%% Tilde in URL in literature
\usepackage{url}
%%% multiple rows
\usepackage{multirow}
%%% fancy column and row seperators
\usepackage{hhline}
%%% custom items in enumerate and itemize
\usepackage{enumitem}

\clubpenalty=5000
\widowpenalty=5000

\setlength{\emergencystretch}{2cm}
%----------------------------------------------------------------------------------------
%	abbreviation includes
%----------------------------------------------------------------------------------------
%%% Standard Math
% todo function
\newcommand{\todo}[1]{\textcolor{red}{#1}}

% spaces
\newcommand{\N}{\ensuremath{\mathbb{N}}}
\newcommand{\R}{\ensuremath{\mathbb{R}}}
\newcommand{\Rpos}{\ensuremath{\mathbb{R}_{>0}}}
\newcommand{\Rnneg}{\ensuremath{\mathbb{R}_{\geq 0}}}
% Complex numbers
%		Problem: \C does not work with XeLaTeX
\newcommand{\Cn}{\ensuremath{\mathbb{C}}}
% imaginary unit
\newcommand{\imagUnit}{\ensuremath{{1i}}}

% formation
\newcommand{\txtfrac}[2]{\ensuremath{{#1}/{#2}}}

% braces
\newcommand{\lb}{\left(}
\newcommand{\rb}{\right)}
\newcommand{\lsb}{\left[}
\newcommand{\rsb}{\right]}
\newcommand{\lcb}{\left\{}
\newcommand{\rcb}{\right\}}

% Operators
\DeclareMathOperator*{\argmin}{\operatornamewithlimits{argmin}}
\DeclareMathOperator*{\argmax}{\operatornamewithlimits{argmax}}
\let\Re\relax
\DeclareMathOperator{\Re}{Re}
\newcommand{\Reb}[1]{\ensuremath{\Re\lb{#1}\rb}}
\let\Im\relax
\DeclareMathOperator{\Im}{Im}
\newcommand{\Imb}[1]{\ensuremath{\Im\lb{#1}\rb}}
\DeclareMathOperator*{\rank}{rank}
\newcommand{\rankb}[1]{\rank\lb{#1}\rb}
\DeclareMathOperator*{\img}{img}
\DeclareMathOperator*{\trace}{trace}
\DeclareMathOperator*{\mspan}{span} % Math SPAN
\DeclareMathOperator*{\colspan}{colspan}
\DeclareMathOperator*{\diag}{diag}
\DeclareMathOperator*{\blkdiag}{blkdiag}
\newcommand{\detb}[1]{\ensuremath{\det{\lb #1 \rb}}}
\newcommand{\blkdiagb}[1]{\ensuremath{\blkdiag{\lb #1 \rb}}}
\newcommand{\diagb}[1]{\ensuremath{\diag{\lb #1 \rb}}}
\newcommand{\dimb}[1]{\ensuremath{\dim{\lb #1 \rb}}}
\newcommand{\expb}[1]{\ensuremath{\exp{\lb #1 \rb}}}
\newcommand{\imgb}[1]{\ensuremath{\img{\lb #1 \rb}}}
\newcommand{\colspanb}[1]{\ensuremath{\colspan{\lb #1 \rb}}}
\newcommand{\traceb}[1]{\trace\lb{#1}\rb}
\newcommand{\spanb}[1]{\ensuremath{\mspan{\lcb #1 \rcb}}}
\newcommand{\sinb}[1]{\ensuremath{\sin{\lb #1 \rb}}}
\newcommand{\abs}[1]{\ensuremath{\left|{#1}\right|}}
\newcommand{\minimize}[1]{\underset{#1}{\operatorname{minimize}}}
\newcommand{\maximize}[1]{\underset{#1}{\operatorname{maximize}}}

%%% Fat alphabet
% Small letters
\newcommand{\fa}{\ensuremath{\bm{a}}}
\newcommand{\fb}{\ensuremath{\bm{b}}}
\newcommand{\fc}{\ensuremath{\bm{c}}}
\newcommand{\fd}{\ensuremath{\bm{d}}}
\newcommand{\fe}{\ensuremath{\bm{e}}}
\newcommand{\ff}{\ensuremath{\bm{f}}}
\newcommand{\fg}{\ensuremath{\bm{g}}}
\newcommand{\fh}{\ensuremath{\bm{h}}}
\newcommand{\fati}{\ensuremath{\bm{i}}}
\newcommand{\fj}{\ensuremath{\bm{j}}}
\newcommand{\fk}{\ensuremath{\bm{k}}}
\newcommand{\fl}{\ensuremath{\bm{l}}}
\newcommand{\fm}{\ensuremath{\bm{m}}}
\newcommand{\fn}{\ensuremath{\bm{n}}}
\newcommand{\fo}{\ensuremath{\bm{o}}}
\newcommand{\fp}{\ensuremath{\bm{p}}}
\newcommand{\fq}{\ensuremath{\bm{q}}}
\newcommand{\fr}{\ensuremath{\bm{r}}}
\newcommand{\fs}{\ensuremath{\bm{s}}}
\newcommand{\ft}{\ensuremath{\bm{t}}}
\newcommand{\fu}{\ensuremath{\bm{u}}}
\newcommand{\fv}{\ensuremath{\bm{v}}}
\newcommand{\fw}{\ensuremath{\bm{w}}}
\newcommand{\fx}{\ensuremath{\bm{x}}}
\newcommand{\fy}{\ensuremath{\bm{y}}}
\newcommand{\fz}{\ensuremath{\bm{z}}}

% Capital letters
\newcommand{\fA}{\ensuremath{\bm{A}}}
\newcommand{\fB}{\ensuremath{\bm{B}}}
\newcommand{\fC}{\ensuremath{\bm{C}}}
\newcommand{\fD}{\ensuremath{\bm{D}}}
\newcommand{\fE}{\ensuremath{\bm{E}}}
\newcommand{\fF}{\ensuremath{\bm{F}}}
\newcommand{\fG}{\ensuremath{\bm{G}}}
\newcommand{\fH}{\ensuremath{\bm{H}}}
\newcommand{\fI}{\ensuremath{\bm{I}}}
\newcommand{\fJ}{\ensuremath{\bm{J}}}
\newcommand{\fK}{\ensuremath{\bm{K}}}
\newcommand{\fL}{\ensuremath{\bm{L}}}
\newcommand{\fM}{\ensuremath{\bm{M}}}
\newcommand{\fN}{\ensuremath{\bm{N}}}
\newcommand{\fO}{\ensuremath{\bm{O}}}
\newcommand{\fP}{\ensuremath{\bm{P}}}
\newcommand{\fQ}{\ensuremath{\bm{Q}}}
\newcommand{\fR}{\ensuremath{\bm{R}}}
\newcommand{\fS}{\ensuremath{\bm{S}}}
\newcommand{\fT}{\ensuremath{\bm{T}}}
\newcommand{\fU}{\ensuremath{\bm{U}}}
\newcommand{\fV}{\ensuremath{\bm{V}}}
\newcommand{\fW}{\ensuremath{\bm{W}}}
\newcommand{\fX}{\ensuremath{\bm{X}}}
\newcommand{\fY}{\ensuremath{\bm{Y}}}
\newcommand{\fZ}{\ensuremath{\bm{Z}}}

% Greek small letters
\newcommand{\falpha}{\ensuremath{\bm{\alpha}}}
\newcommand{\fbeta}{\ensuremath{\bm{\beta}}}
\newcommand{\fgamma}{\ensuremath{\bm{\gamma}}}
\newcommand{\fdelta}{\ensuremath{\bm{\delta}}}
\newcommand{\fepsilon}{\ensuremath{\bm{\vareps}}}
\newcommand{\fzeta}{\ensuremath{\bm{\zeta}}}
\newcommand{\feta}{\ensuremath{\bm{\eta}}}
\newcommand{\ftheta}{\ensuremath{\bm{\theta}}}
\newcommand{\flota}{\ensuremath{\bm{\lota}}}
\newcommand{\fkappa}{\ensuremath{\bm{\kappa}}}
\newcommand{\flambda}{\ensuremath{\bm{\lambda}}}
\newcommand{\fmu}{\ensuremath{\bm{\mu}}}
\newcommand{\fnu}{\ensuremath{\bm{\nu}}}
\newcommand{\fxi}{\ensuremath{\bm{\xi}}}
% skipped omikron
\newcommand{\fpi}{\ensuremath{\bm{\pi}}}
\newcommand{\frho}{\ensuremath{\bm{\rho}}}
\newcommand{\fsigma}{\ensuremath{\bm{\sigma}}}
\newcommand{\ftau}{\ensuremath{\bm{\tau}}}
\newcommand{\fupsilon}{\ensuremath{\bm{\upsilon}}}
\newcommand{\fphi}{\ensuremath{\bm{\varphi}}}
\newcommand{\fchi}{\ensuremath{\bm{\chi}}}
\newcommand{\fpsi}{\ensuremath{\bm{\psi}}}
\newcommand{\fomega}{\ensuremath{\bm{\omega}}}

% some Greek capital letters
\newcommand{\fGamma}{\ensuremath{\bm{\Gamma}}}
\newcommand{\fDelta}{\ensuremath{\bm{\Delta}}}
\newcommand{\fTheta}{\ensuremath{\bm{\Theta}}}
\newcommand{\fLambda}{\ensuremath{\bm{\Lambda}}}
\newcommand{\fXi}{\ensuremath{\bm{\Xi}}}
\newcommand{\fPi}{\ensuremath{\bm{\Pi}}}
\newcommand{\fSigma}{\ensuremath{\bm{\Sigma}}}
\newcommand{\fUpsilon}{\ensuremath{\bm{\Upsilon}}}
\newcommand{\fPhi}{\ensuremath{\bm{\varPhi}}}
\newcommand{\fPsi}{\ensuremath{\bm{\Psi}}}
\newcommand{\fOmega}{\ensuremath{\bm{\varOmega}}}

% Numbers
\newcommand{\fzero}{\ensuremath{\bm{0}}}
\newcommand{\fone}{\ensuremath{\bm{1}}}

%%% Standard mathematical operators
% (real) Transpose operator
%		Problem: \T does not work with XeLaTeX
\newcommand{\rT}[1]{\ensuremath{#1^{\textsf{T}}}}
\newcommand{\rTb}[1]{\ensuremath{\rT{\lb#1\rb}}}
\newcommand{\rTsb}[1]{\ensuremath{\rT{\lsb#1\rsb}}}

% Conjugate transpose
\newcommand{\cT}[1]{\ensuremath{#1^{\ast}}}
\newcommand{\cTb}[1]{\ensuremath{\cT{\lb{}#1\rb}}}

% Hermitian operator
\newcommand{\Her}[1]{\ensuremath{{#1}^{\textsf{H}}}}
\newcommand{\Herb}[1]{\ensuremath{\Her{\lb{}#1\rb}}}

% Symplectic inverse
\newcommand{\si}[1]{\ensuremath{#1^{+}}}
\newcommand{\sib}[1]{\ensuremath{\si{\lb{}#1\rb}}}

% orthogonal complement
\newcommand{\orth}[1]{\ensuremath{#1^{\perp}}}

% Inverse operator
\newcommand{\inv}[1]{\ensuremath{{#1}^{\textsf{-1}}}}
\newcommand{\invb}[1]{\ensuremath{\inv{\lb{#1}\rb}}}

% Inverse transposed operator
\newcommand{\invT}[1]{\ensuremath{{#1}^{\textsf{-T}}}}
\newcommand{\invTb}[1]{\ensuremath{\invT{\lb{#1}\rb}}}

% Gradients
\newcommand{\grad}[1][]{\ensuremath{\nabla_{#1}}}

% Divergence
\let\div\relax
\DeclareMathOperator{\div}{div}
\newcommand{\divb}[2][]{\ensuremath{\div_{#1}\lb{#2}\rb}}

% parial derivative
\newcommand{\deldel}[2][]{\ensuremath{\frac{\partial #1}{\partial #2}}}
\newcommand{\deldelt}{\ensuremath{\deldel{t}}}
\newcommand{\deldelfx}{\ensuremath{\deldel{\fx}}}

% second-order partial derivative
\newcommand{\sodeldel}[1]{\ensuremath{\frac{\partial^2}{\partial^2 #1}}} % Second-Order DELDEL

% jacobian
\newcommand{\jac}[2]{\deldel[#1]{#2}}
\newcommand{\jacp}[2]{\lb
    \jac{#1}{#2}
\rb}

% Derivatives written out
\newcommand{\dd}[2][]{\ensuremath{\frac{\mathrm{d}^{#1}}{\mathrm{d}#2}}}
\newcommand{\ddfx}{\ensuremath{\dd{\fx}}}
\newcommand{\ddfy}{\ensuremath{\dd{\fy}}}
\newcommand{\ddt}{\ensuremath{\dd{t}}}
\newcommand{\dds}{\ensuremath{\dd{s}}}

% Second order derivative written out
\newcommand{\soddt}{\dd[2]{t^2}}

% Integral operators
\newcommand{\intd}[1]{\ensuremath{\mathrm{d}#1}} % integral d (\d already in use)
\newcommand{\ds}{\ensuremath{\intd{s}}}

% Symplectic matrix
\newcommand{\Jtwo}[1]{\ensuremath{\mathbb{J}_{2#1}}}
\newcommand{\TJtwo}[1]{\ensuremath{\rT{\mathbb{J}}_{2#1}}}
\newcommand{\Jtn}{\ensuremath{\Jtwo{n}}}
\newcommand{\TJtn}{\ensuremath{\TJtwo{n}}}
\newcommand{\Jtk}{\ensuremath{\Jtwo{k}}}
\newcommand{\TJtk}{\ensuremath{\TJtwo{k}}}
\newcommand{\J}{\ensuremath{\mathbb{J}}}

% Identity matrix
\newcommand{\I}[1]{\ensuremath{\fI_{#1}}}
\newcommand{\In}{\ensuremath{\I{n}}}

% Matrix of zeros
\newcommand{\Z}[1]{\ensuremath{\fzero}_{#1}}
\newcommand{\Zn}{\ensuremath{\Z{n}}}

% Matrix of ones
\newcommand{\Ones}[1]{\ensuremath{\fone}_{#1}}

% Norms
\newcommand{\norm}[1]{\left\lVert{#1}\right\rVert}
\newcommand{\wnorm}[2]{\ensuremath{\norm{#2}_{#1}}} % weighted norm
\newcommand{\Fnorm}[1]{\wnorm{\mathrm{F}}{#1}}
\newcommand{\tnorm}[1]{\wnorm{\mathrm{2}}{#1}} % two norm

% Inner products
\newcommand{\ip}[2]{\left\langle #1,\; #2 \right\rangle}
\newcommand{\wip}[3]{\ip{#2}{#3}_{#1}}

% Special matrices
\newcommand{\sq}[1]{\ensuremath{{#1}^{\frac{1}{2}}}}
\newcommand{\isq}[1]{\ensuremath{{#1}^{-\frac{1}{2}}}}

% standard ODE stuff
\newcommand{\tInit}{\ensuremath{t_{\mathrm{0}}}}
\newcommand{\tEnd}{\ensuremath{t_{\mathrm{end}}}}
\newcommand{\ftInterval}{\ensuremath{[\tInit, \tEnd]}} % finite time interval
\newcommand{\fxInit}{\ensuremath{\fx_{\mathrm{0}}}}

% thesis specific commands

\newcommand{\qpvec}{\begin{pmatrix}
    q \\
    p
\end{pmatrix}}

\newcommand{\uppersympop}[1]{\begin{bmatrix}
    id & {#1} \\
    0 & id
\end{bmatrix} \qpvec}

\newcommand{\lowersympop}[1]{\begin{bmatrix}
    id & 0 \\
    {#1} & id
\end{bmatrix} \qpvec}

\newcommand{\activation}{\sigma}

\newcommand{\onevec}[1]{1_{#1}}

\newcommand*\phantomrel[1]{\mathrel{\phantom{#1}}}

\newcommand{\up}{\text{up}}
\newcommand{\low}{\text{low}}
\newcommand{\layertf}{\mathrm{T}}
%%% Clever refing
\usepackage[capitalise,noabbrev]{cleveref}

%----------------------------------------------------------------------------------------
%	References with cleveref
%----------------------------------------------------------------------------------------
\crefformat{equation}{(#2#1#3)}

%----------------------------------------------------------------------------------------
%	Theorem environments
%----------------------------------------------------------------------------------------
\newtheorem{theorem}{Theorem}
\newtheorem{example}{Example}
\newtheorem{corollary}{Corollary}
\newtheorem{lemma}{Lemma}

%----------------------------------------------------------------------------------------
%	fancyhdr
%----------------------------------------------------------------------------------------
\usepackage{fancyhdr}

\makeatletter
\newcommand{\theauthor}{Steffen Müller} %
\makeatother

\renewcommand{\headrulewidth}{0.3pt}
\renewcommand{\footrulewidth}{0.3pt}
\fancyfoot[OL,ER]{University of Stuttgart} % inner
\fancyfoot[OR,EL]{\thepage} % outer
\fancyhead[OL,ER]{\theauthor} % inner 
\fancyhead[OR,EL]{IANS -- Institute of Applied Analysis and Numerical Simulation} % outer
\fancyfoot[C]{}

% \headsep=4mm
% \footskip=4mm
\parindent=0mm
\parskip=6pt
% \renewcommand{\footskip}{3pt}

%----------------------------------------------------------------------------------------
%	Something
%----------------------------------------------------------------------------------------
\usepackage{textpos}
\setlength{\TPHorizModule}{1mm}%
\setlength{\TPVertModule}{1mm}%
% \headsep=5mm
% \footskip=5mm
\pagestyle{fancy}
\newcommand{\articleheading}[3]{
{\large #1}\\[3mm]
{\Large\bf #2}\\[3mm]
{\large #3}
}

%----------------------------------------------------------------------------------------
%	Start Document
%----------------------------------------------------------------------------------------
\begin{document}
\pagenumbering{roman}
%----------------------------------------------------------------------------------------
%
% TITLE PAGE
%
%----------------------------------------------------------------------------------------
%------------------------------------------
\begin{titlePageST}
%------------------------------------------
\makeLogo%
{-10pt}{
\includegraphics[width=0.7\textwidth]{figures/logos/simtech.pdf}
}%
{0pt}{
	\begin{center}
		\includegraphics[width=0.4\textwidth]{figures/logos/ians.pdf}
	\end{center}}%
{0pt}{
\begin{flushright}
	\vspace{-10pt}
	\includegraphics[width=0.9\textwidth]{figures/logos/unistuttgart_logo_englisch_cmyk.eps}
\end{flushright}}%
\vspace{35pt}%
%------------------------------------------
\makeHeader%
[Research Group: Numerical Mathematics] %
{Institute of Applied Analysis and Numerical Simulation} %
\vspace{50pt}%
%------------------------------------------
\makeTitle%
{Simulation Technology Degree Course} %
{Bachelor Thesis} %
\vspace{80pt}%
%------------------------------------------
\makeTitleThesis%
{Symplectic Neural Networks}
\vspace{90pt}%
%------------------------------------------
\begin{supervisorST}{3}%
\addSuper%
{First Reviewer}%
{Prof. Dr. B. Haasdonk}%
{Institute of Applied Analysis and\\[-0.2cm]
Numerical Simulation (IANS)}%
\addSuper%
{Second Reviewer}%
{Prof. Dr. D. Pflüger}%
{Institute of Parallel and Distributed\\[-0.2cm]
Systems (Scientific Computing)}%
\addSuper%
{Advisor}%
{Patrick Buchfink, M.Sc.}%
{Institute of Applied Analysis and\\[-0.2cm]
Numerical Simulation (IANS)}%
\end{supervisorST}%
\vspace{80pt}%
%------------------------------------------
\begin{authorST}{Submitted by}%
\addAuthorInfo{Author}{Steffen Müller}
\addAuthorInfo{Student ID}{3260643}
\addAuthorInfo{SimTech ID}{119} %
\addAuthorInfo{Submission Date}{...} %FILLIN
\end{authorST}%
%------------------------------------------
\end{titlePageST}
%----------------------------------------------------------------------------------------
\clearpage
\newpage\thispagestyle{plain}\null
\newpage\thispagestyle{plain}

%----------------------------------------------------------------------------------------
%
% ABSTRACT
%
%----------------------------------------------------------------------------------------
\section*{Abstract}
...
\clearpage
\newpage\thispagestyle{plain}\null
\newpage\thispagestyle{plain}

%----------------------------------------------------------------------------------------
%
% ACKNOWLEDGEMENTS
%
%----------------------------------------------------------------------------------------
\section*{Acknowledgements}
...
\clearpage
\newpage\thispagestyle{plain}\null
\newpage\thispagestyle{plain}

%----------------------------------------------------------------------------------------
%
% TOC
%
%----------------------------------------------------------------------------------------
\setcounter{tocdepth}{2}
\tableofcontents
\clearpage
\newpage\thispagestyle{plain}\null
%----------------------------------------------------------------------------------------
%
% Begin with document content
%
%----------------------------------------------------------------------------------------
\newpage
\pagenumbering{arabic} 
%----------------------------------------------------------------------------------------
%
% Introduction
%
%----------------------------------------------------------------------------------------
\section{Introduction}

\todo{Introduce Hamiltonain systems with general skew-symmetric $S$?}

\subsection{Outline}

\subsection{Notation}

Let $f: \mathbb{R} \rightarrow \mathbb{R}$ be an arbitrary function and $x \in \mathbb{R}^n$
a vector. When we write $f(x)$ we mean the element-wise application of $f$ 
on the vector $x \in \mathbb{R}^n$, i.e. $f(x) = \lb
	f(x_1), ..., f(x_n)
\rb^T$.

The $d$-by-$d$ identity matrix is denoted by $I_d$. If the dimension $d$ can be inferred
from context we may just write $I$.

$\onevec{n}$ denotes the $n$-dimensional $1$-vector.

\todo{Introduce block matrix notation for general functions}

We denote the $k$-th partial derivative of a function $f: \mathbb{R}^n \to \mathbb{R},
x \mapsto f(x)$ with $\deldel{x_k} f(x)$.

If $f: \mathbb{R}^{n_1} \to \mathbb{R}^{n_2}$ is a multi-dimensional function

\todo{"partial derivative" of a vector, Jacobian matrices}

\todo{Gradient}

\todo{Describe $d$}

%----------------------------------------------------------------------------------------
%
% Content
%
%----------------------------------------------------------------------------------------
\newpage
\section{Problem setup}

\section{Architecture}

In this section we describe the architecture of our symplectic neural networks.
The linear layers, activation layers and gradient layers are recapitulated as 
proposed by \citeauthor{Jin2020} in \cite{Jin2020}.
The convolution layers and normalized variants of the activation and gradient layers
are contributions of this work. 
\todo{Convolution gradient module, refer to similar reference}
We supplement proofs for symplecticity for all layer types. 
A similar architecture was initially proposed by \citeauthor{Deco1995} in \cite{Deco1995} 
\todo{Have a deeper look at \cite{Deco1995}}.

\todo{we use that the composition of symplectic maps is again symplectic}

\todo{Relation to Residual Networks / ResNet}

\begin{lemma}\label{jacobi_symmetric}
	Let
	\begin{equation*}
		f_{up} \qpvec := \uppersympop{g}
	\end{equation*}
	and
	\begin{equation*}
		f_{low} \qpvec := \lowersympop{g}
	\end{equation*}
	where $g: \mathbb{R}^d \to \mathbb{R}^d,\; p \mapsto g(p)$ is a function in 
	$C^1(\mathbb{R}^d, \mathbb{R}^d)$. 

	$f_{up}$ and $f_{low}$ are symplectic if and only if the Jacobian matrix $\jac{g}{p}$
	is symmetric.
\end{lemma}
\begin{proof}
	We show the result for $f_{up}$ only. The proof is analogous for $f_{low}$.

	\begin{equation*}
		\deldel[f_{up}]{(q,p)} = \begin{pmatrix}
			I & \deldel[g]{p} \\
			0 & I
		\end{pmatrix}
	\end{equation*}
	It follows
	\begin{align*}
		\jacp{f_{up}}{(q,p)}^T J \jacp{f_{up}}{(q,p)} 
		&= \left(\deldel[f_{up}]{(q,p)}\right)^T \begin{pmatrix}
			0 & I \\
			-I & -\jac{g}{p}
		\end{pmatrix} \\
		&= \begin{pmatrix}
			I & 0 \\
			\jacp{g}{p}^T & Id
		\end{pmatrix} \begin{pmatrix}
			0 & I \\
			-I & -\jac{g}{p}
		\end{pmatrix} \\
		&= \begin{pmatrix}
			0 & I \\
			-I & \jacp{g}{p}^T-\jac{g}{p}
		\end{pmatrix}
	\end{align*}
	Thus $f_{up}$ is symplectic if and only if $\jacp{g}{p}^T-\jac{g}{p}=0$, 
	i.e. if and only if the Jacobian $\deldel[g]{p}$ is everywhere symmetric.
\end{proof}

\begin{corollary}\label{matrix_symmetric}
	Let
	\begin{equation*}
		f_{up} \qpvec := \begin{pmatrix}
			I & S \\
			0 & I
		\end{pmatrix} \qpvec
	\end{equation*}
	and
	\begin{equation*}
		f_{low} \qpvec := \begin{pmatrix}
			I & 0 \\
			S & I
		\end{pmatrix} \qpvec
	\end{equation*}
	where $S \in \mathbb{R}^{dxd}$. 

	$f_{up}$ and $f_{low}$ are symplectic if and only if the matrix $S$
	is symmetric.
\end{corollary}

The next corollary is useful to construct symplectic maps.
\begin{corollary}\label{gradient_corollary}
	Let $V: \mathbb{R}^d \to \mathbb{R}, \; p \mapsto V(p)$ be a function in 
	$C^2(\mathbb{R}^d, \mathbb{R})$. 
	
	Then
	\begin{equation*}
		f_{up} \qpvec := \uppersympop{\grad{V}}
	\end{equation*}
	and
	\begin{equation*}
		f_{low} \qpvec := \lowersympop{\grad{V}}
	\end{equation*}
	define a symplectic map. We call $V$ a potential.
\end{corollary}
\begin{proof}
	The Jacobian matrix of $\grad{V}$ corresponds to the Hessian matrix of $V$,
	i.e. $\jac{(\grad{V})}{p} = HD$. \todo{Notation for Hessian}
	The Hessian is symmetric, thus the result follows with \cref{jacobi_symmetric}.
\end{proof}

\todo{Relation to 'Generating functions' in symplectic literature?}

\subsection{Linear layers}

\begin{align*}
	\ell_{up}: \mathbb{R}^{2d} \to \mathbb{R}^{2d}, \quad \ell_{up} \qpvec &:= \begin{pmatrix}
		I & S \\
		0 & I
	\end{pmatrix} \qpvec + b, \\
	\ell_{low}: \mathbb{R}^{2d} \to \mathbb{R}^{2d}, \quad \ell_{low} \qpvec &:= \begin{pmatrix}
		I & 0 \\
		S & I
	\end{pmatrix} \qpvec + b
\end{align*}
where $S \in \mathbb{R}^{d \times d}$ symmetric and $b \in \mathbb{R}^{2d}$
are learnable parameters. In practice, we parametrize the symmetric matrix $S\in \mathbb{R}^{d \times d}$
with $S = A^T + A$ via another arbitrary matrix $A\in \mathbb{R}^{d \times d}$, as
most optimization methods used for learning neural networks are designed for
unconstrained optimization problems.

It follows directly from \cref{matrix_symmetric} that the linear layers are
symplectic. Alternatively, we can apply \cref{gradient_corollary} by
choosing the potential $V(p) := p^TAp$.

\begin{proof}
	Let
	\begin{align*}
		\phi_1: \mathbb{R}^d \to \mathbb{R}^{2d}, \quad &\phi_1(p) := \begin{pmatrix}
			p \\
			Ap
		\end{pmatrix} \\
		\phi_2: \mathbb{R}^{2d} \to \mathbb{R}, \quad &\phi_2(p, \hat{p}) := 
		p^T\hat{p}
	\end{align*}
	
	Then $V(p) = \phi_2(\phi_1(p))$ and the corresponding Jacobian matrices are
	\begin{align*}
		\jac{\phi_1}{p}(p) &= \begin{pmatrix}
			I \\
			A
		\end{pmatrix} \\
		\jac{\phi_2}{(p,\hat{p})}(p, \hat{p}) &= \begin{pmatrix}
			\hat{p}^T && p^T
		\end{pmatrix}
	\end{align*}
	Thus
	\begin{align*}
		\jac{V}{p}(p) &= 
		\jac{\phi_2}{(p,\hat{p})} (\phi_1(p))
		\jac{\phi_1}{p}(p) \\
		&= \begin{pmatrix}
			(Ap)^T && p^T
		\end{pmatrix}
		\begin{pmatrix}
			I \\
			A
		\end{pmatrix} \\
		&= (Ap)^T + p^TA = (Ap + A^Tp)^T \\
		&= ((A+A^T)p)^T
	\end{align*}
	Therefore $\grad{V}(p) = \lb \deldel{p}V(p) \rb^T = (A+A^T)p = Sp$ with
	$S := A+A^T$ symmetric. Symplecticity follows with \cref{gradient_corollary}.
\end{proof}

We may enhance expressivity of a linear layer by alternately composing multiple
$\ell_{up}$ and $\ell_{low}$:

\begin{align*}
	\mathcal{L}^{n}_{up} \qpvec &= \begin{pmatrix}
		I && 0 / S_n \\
		S_n / 0 && I
	\end{pmatrix}
	\cdots
	\begin{pmatrix}
		I && 0 \\
		S_2 && I
	\end{pmatrix}
	\begin{pmatrix}
		I && S_1 \\
		0 && I
	\end{pmatrix} + b \\
	\mathcal{L}^{n}_{low} \qpvec &= \begin{pmatrix}
		I && S_n / 0 \\
		0 / S_n && I
	\end{pmatrix}
	\cdots
	\begin{pmatrix}
		I && S_2 \\
		0 && I
	\end{pmatrix}
	\begin{pmatrix}
		I && 0 \\
		S_1 && I
	\end{pmatrix} + b
\end{align*}

$\mathcal{L}^{n}_{up}$ and $\mathcal{L}^{n}_{low}$ are again symplectic because they 
are a composition of the symplectic maps $\ell_{up}$ and $\ell_{low}$.

\citeauthor{jin2020unit} show in \cite{jin2020unit} that $\mathcal{L}^{9}_{up}$
can parametrize every symplectic linear map. In other words, 
the set of all possible $\mathcal{L}^{9}_{up}$ is equal to the set of all symplectic linear maps.
\todo{check how \cite{jin2020unit} appears in bibliography}

It does not make sense to put two upper linear layers or two lower linear layers after each other,
because this reduces to a single upper or lower linear layer:

\begin{proof}
	We show the statement for two upper linear layers.
	\begin{align*}
		\ell_{up,2}\lb\ell_{up,1} \qpvec\rb &=
		\begin{pmatrix}
			I && S_2 \\
			0 && I
		\end{pmatrix}
		\lb
		\begin{pmatrix}
			I && S_1 \\
			0 && I
		\end{pmatrix}
		\qpvec + b_1
		\rb + b_2 \\
		&= \begin{pmatrix}
			I && S_1 + S_2 \\
			0 && I
		\end{pmatrix} \qpvec
		+ \begin{pmatrix}
			I && S_2 \\
			0 && I
		\end{pmatrix} b_1
		+ b_2 \\
		&= \begin{pmatrix}
			I && S \\
			0 && I
		\end{pmatrix} \qpvec + b
	\end{align*}
	with $S := S_1 + S_2$ and $b := \begin{pmatrix}
		I && S_2 \\
		0 && I
	\end{pmatrix} b_1
	+ b_2$.
\end{proof}

\subsection{Convolution layers}

\todo{Try motivating convolution layers by showing that the linear layers explode if $d$ becomes large.}

Given a kernel $k \in \mathbb{R}^{s_k}$ with kernel size $s_k \in \mathbb{N}$
the $1$-dimensional convolution operator is defined as
\begin{equation*}
	conv1d[k]: \mathbb{R}^{L_{in}} \to \mathbb{R}^{L_{out}}, \quad
	\lb conv1d[k](x) \rb_i := \sum_{u=1}^{s_k} k_u x_{i+u-1} \quad
	(1 \leq i \leq L_{out})
\end{equation*}
where $L_{in}$ is the input dimension and $L_{out} = L_{in} - s_k + 1$.
\todo{Briefly mention more general versions with stride and dilation.}
The convolution operator reduces the input dimension $L_{in}$ (if $s_k > 1$).

\todo{Add illustration for convolution}

As we will justify later we only use kernels with odd kernel size $s_k = 2m+1$, $m \in \mathbb{N}$.
Therefore from now on we always implicitly assume a kernel with odd kernel size $s_k = 2m+1$.
\todo{in this section, later conv gradient module not?}
To simplify notation we introduce new indexing for kernels by counting from $-m$ to $m$,
i.e. $k = (k_{-m}, k_{-m+1}, \dots, k_{-1} k_0, k_1, \dots, k_{m})^T \in \mathbb{R}^{2m+1}$.
With new indexing the convolution operator becomes 
\begin{equation*}
	\lb conv1d[k](x) \rb_i = \sum_{u=-m}^{m} k_u x_{i+u+m} \quad
	(1 \leq i \leq L_{out})
\end{equation*}

We call a kernel $k \in \mathbb{R}^{2m+1}$ symmetric if $k_i = k_{-i}$ for all $i = 0, \dots, m$.

A padding operator increases the input dimension by padding
the input boundaries. Let us define $1$-dimensional constant padding and symmetric padding operators.
Let $s_p \in \mathbb{N}$ be the padding size and $L_{in}$ the input dimension.

Constant padding adds constant padding values at the left and right boundary.
\begin{align*}
	&constant\_pad[s_p,l,r] : \mathbb{R}^{L_{in}} \to \mathbb{R}^{L_{in}+2s_p}, \\[7pt]
	&(constant\_pad[s_p,l,r](x))_i :=
	\begin{cases}
		l_i & 1 \leq i \leq s_p \\
		x_{i-s_p} & s_p+1 \leq i \leq L_{in}+s_p \\
		r_{i-L_{in}+s_p} & L_{in}+s_p+1 \leq i \leq L_{in} + 2s_p
	\end{cases}
	\quad (1 \leq i \leq L_{in}+2s_p)
\end{align*}
where $l,r \in \mathbb{R}^{s_p}$ are the padding values for the left and right boundaries.

Symmetric padding mirrors the inner values at the left and right boundaries.
\begin{align*}
	&symmetric\_pad[s_p] : \mathbb{R}^{L_{in}} \to \mathbb{R}^{L_{in}+2s_p}, \\[7pt]
	&(symmetric\_pad[s_p](x))_i :=
	\begin{cases}
		x_{s_p-i+1} & 1 \leq i \leq s_p \\
		x_{i-s_p} & s_p+1 \leq i \leq L_{in}+s_p \\
		x_{2L_{in}+s_p+1-i} & L_{in}+s_p+1 \leq i \leq L_{in} + 2s_p
	\end{cases}
	\quad (1 \leq i \leq L_{in}+2s_p)
\end{align*}
For symmetric padding it must hold $L_{in} \geq s_p$.

\todo{Add illustrations for padding}

We define the symplectic $1$-dimensional convolution layer by
\begin{align*}
	\mathcal{C}_{up} : \mathbb{R}^{2d} \to \mathbb{R}^{2d},
	\mathcal{C}_{up} \qpvec &:= \uppersympop{conv1d[k] \circ pad[s_p,\dots]} \\
	\mathcal{C}_{low} : \mathbb{R}^{2d} \to \mathbb{R}^{2d},
	\mathcal{C}_{low} \qpvec &:= \lowersympop{conv1d[k] \circ pad[s_p,\dots]}
\end{align*}
where kernel $k \in \mathbb{R}^{2m+1}$ symmetric, $s_p = m$ and either
\begin{align*}
	pad[s_p,\dots] &= constant\_pad[s_p,l,r],\; l,r \in \mathbb{R}^{s_p} \text{ or } \\
	pad[s_p,\dots] &= symmetric\_pad[s_p]
\end{align*}
For $pad[s_p,\dots] = symmetric\_pad[s_p]$ it must hold $d \geq m$.
The kernel $k \in \mathbb{R}^{2m+1}$ is learnable during training.
If using constant pad $l,r \in \mathbb{R}^{s_p}$ can be either fixed or learnable during training.

Indeed $conv1d[k] \circ pad[p,\dots]$ maps from $\mathbb{R}^d$ to $\mathbb{R}^d$ because
\begin{align*}
	L_{in,pad[s_p,\dots]} &= d \\
	L_{out,pad[s_p,\dots]} &= L_{in,pad[s_p,\dots]} + 2s_p = d+2s_p = d+2m \\
	L_{in,conv1d[k]} &= L_{out,pad[p,\dots]} = d+2m \\
	L_{out,conv1d[k]} &= L_{in,conv1d[k]} - s_k + 1 = (d+2m) - (2m+1) + 1 = d
\end{align*}

Furthermore the two conditions
\begin{align}
	L_{in,conv1d[k]} &= L_{out,pad[s_p,\dots]} = d+2s_p \label{conv_dim1} \\ 
	L_{out,conv1d[k]} &= L_{in,conv1d[k]} - s_k + 1 = d \label{conv_dim2}
\end{align}
indeed imply that the kernel size $s_k$ must be odd.
\begin{equation*}
	\text{\cref{conv_dim1} in \cref{conv_dim2}:} \quad d+2s_p - s_k + 1 = d \iff s_k = 2s_p+1
\end{equation*}

Let ${b_1, b_2, \dots b_m} \in \mathbb{R}^{2m+1}$ be a basis of the space 
$\{ k \in \mathbb{R}^{2m+1} : k \text{ is symmetric} \}$.
The symmetric kernel $k$ is then parametrized by the coefficients $\beta \in \mathbb{R}^m$ via
\begin{equation*}
	k = (b_1, b_2, \dots, b_m) \beta = \beta_1 b_1 + \beta_2 b_2 + \dots + \beta_m b_m
\end{equation*}

We can choose the canonical basis with basis vectors
\begin{align*}
	b_1^C = (0, \dots, 0, 0, 0, &1, 0, 0, 0, \dots, 0)^T \in \mathbb{R}^{2m+1}, \\
	b_2^C = (0, \dots, 0, 0, 1, &0, 1, 0, 0, \dots, 0)^T \in \mathbb{R}^{2m+1}, \\
	b_3^C = (0, \dots, 0, 1, 0, &0, 0, 1, 0, \dots, 0)^T \in \mathbb{R}^{2m+1}, \\
	&\vdots \\
	b_m^C = (1, 0, \dots, &0, \dots, 0, 1)^T \in \mathbb{R}^{2m+1}
\end{align*}

Another possible basis choice inspired by finite differences is
\todo{Cite paper with similar idea?}
\begin{align*}
	b_1^{FD} = (0, \dots, 0,& 1,0, \dots, 0)^T \in \mathbb{R}^{2m+1},\\
	b_2^{FD} = (0, \dots, 0, 1,-&2,1, 0, \dots, 0)^T \in \mathbb{R}^{2m+1},\\
	b_3^{FD} = (0, 0, \dots, 0, 1,-4,& 6,-4,1, 0, \dots, 0)^T \in \mathbb{R}^{2m+1} \\
	&\vdots
\end{align*}
The entries for a basis vector $b_i \in \mathbb{R}^{2m+1}$ $(1 \leq i \leq m)$ 
originate from Pascal's triangle.
\begin{equation*}
	\lb b_i^{FD} \rb_j = \begin{dcases}
		(-1)^j \binom{2(i-1)}{j+i-1} : \abs{j} < i \\
		0 : else
	\end{dcases}
	\quad (1 \leq i \leq m, -m \leq j \leq m )
\end{equation*}
where $j$ is indexed from $-m$ to $m$ in accordance to how we defined indexing for kernel entries.

\todo{Move this to Appendix?}
The symmetry of the kernel basis vectors $b_i$ follows with the definition of the binomial coefficient.
\begin{align*}
	\binom{2(i-1)}{j+i-1} &= \frac{2(i-1)!}{(j+i-1)!(2(i-1)-(j+i-1))!} \\
	&= \frac{2(i-1)!}{(j+i-1)!(i-j-1)!}
\end{align*}
Thus we have
\begin{equation*}
	\binom{2(i-1)}{(-j)+i-1} = \binom{2(i-1)}{j+i-1}
\end{equation*}
and
\begin{align*}
	\lb b_i^{FD} \rb_{-j} &= \begin{dcases}
		(-1)^{-j} \binom{2(i-1)}{(-j)+i-1} : \abs{-j} < i \\
		0 : else
	\end{dcases} \\
	&= \begin{dcases}
		(-1)^{j} \binom{2(i-1)}{j+i-1} : \abs{j} < i \\
		0 : else
	\end{dcases} = (b_i^{FD})_j
\end{align*}
for $i=1, \dots, m$ and $j= 0, \dots, m$. So $b_i^{FD}$ $(i=1, \dots, m)$ are a valid basis of
the space $\{ k \in \mathbb{R}^{2m+1} : k \text{ is symmetric} \}$.

\todo{Interesting: Finite difference basis only contains even-order derivatives, otherwise it
would not be symplectic (symmetric). What happens if Hamiltonian contains even-order derivative?
Symmetric Hessian does not explain it, because we only take gradient of Hamiltonian and not Hessian.
Relation to Störmer-Verlet.}

As it turns out in our numerical experiments
parametrization plays an important role how well a neural network learns.

\todo{Maybe move the proof of these theorems to Appendices to keep this cleaner?}

\begin{theorem}
	The $1$-dimensional convolution layers $\mathcal{C}_{up}$ and $\mathcal{C}_{low}$
	with $pad[m,\dots] = constant\_pad[m,l,r]$ ($l,r \in \mathbb{R}^p$) are symplectic.
\end{theorem}
\begin{proof}
	We show that the Jacobian matrix 
	\begin{equation*}
		\jac{\lb conv1d \lsb k \rsb \circ constant\_pad \lsb m,l,r \rsb \rb}{p}
	\end{equation*}
	is symmetric according to \cref{jacobi_symmetric}.

	We have
	\begin{equation*}
		\lb \jac{conv1d \lsb k \rsb}{x} \rb_{ij} = \begin{cases}
			k_{j-i-m} & : \abs{j-i-m} \leq m \\
			0 & : else
		\end{cases}
		\quad (1 \leq i \leq d, 1 \leq j \leq d+2m)
	\end{equation*}
	and
	\begin{equation*}
		\lb \jac{constant\_pad \lsb m,l,r \rsb}{x} \rb_{ij} = \begin{cases}
			1 & : i-j=m \\
			0 & : else
		\end{cases}
		\quad (1 \leq i \leq d+2m, 1 \leq j \leq d)
	\end{equation*}
	Let $1 \leq i,j \leq d$. By the chain rule we have
	\begin{align*}
		\lb \jac{\lb conv1d \lsb k \rsb \circ constant\_pad \lsb m,l,r \rsb \rb}{p} \rb_{ij}
		&= \lb  
			\jac{conv1d \lsb k \rsb}{x}
			\jac{constant\_pad \lsb m,l,r \rsb}{x} 
		\rb_{ij} \\
		&= \sum_{u=1}^{d+2m} \lb \jac{conv1d \lsb k \rsb}{x} \rb_{iu}
		\lb \jac{constant\_pad \lsb m,l,r \rsb}{x} \rb_{uj} \\
		&= \lb \jac{conv1d \lsb k \rsb}{x} \rb_{i(j+m)}
		\quad (\text{bc. } u-j=m \iff u=j+m) \\
		&= \begin{cases}
			k_{(j+m)-i-m} & : \abs{(j+m)-i-m} \leq m \\
			0 & : else
		\end{cases} \\
		&= \begin{cases}
			k_{j-i} & : \abs{j-i} \leq m \\
			0 & : else
		\end{cases}
	\end{align*}
	Then 
	\begin{align*}
		\lb \jac{\lb conv1d \lsb k \rsb \circ constant\_pad \lsb m,l,r \rsb \rb}{p} \rb_{ji}
		&= \begin{cases}
			k_{i-j} & : \abs{i-j} \leq m \\
			0 & : else
		\end{cases} \\
		&= \begin{cases}
			k_{j-i} & : \abs{j-i} \leq m \\
			0 & : else
		\end{cases} \quad \text{($k$ symmetric)} \\
		&= \lb \jac{\lb conv1d \lsb k \rsb \circ constant\_pad \lsb m,l,r \rsb \rb}{p} \rb_{ij}
	\end{align*}
	We have shown that the Jacobian matrix is symmetric.
\end{proof}

\begin{theorem}
	The $1$-dimensional convolution layers $\mathcal{C}_{up}$ and $\mathcal{C}_{low}$
	with $pad[m,\dots] = symmetric\_pad[m]$ are symplectic.
\end{theorem}
\begin{proof}
	Let $1 \leq i \leq d$.
	Then
	\begin{equation*}
		\lb conv1d \lsb k \rsb \circ symmetric\_pad \lsb m \rsb(p) \rb_i
		= \sum_{u=-m}^m k_u \lb symmetric\_pad \lsb m \rsb(p) \rb_{i+u+m} = (*)
	\end{equation*}
	With the variable transform $v=v(u):=i+u+m \iff u=v-i-m$ we have
	\begin{equation*}
		(*) = \sum_{v=i}^{i+2m} k_{v-i-m} \lb symmetric\_pad \lsb m \rsb(p) \rb_{v}
	\end{equation*}
	Let $I := \{ v \in \mathbb{N} : i \leq v \leq i+2m \}$ and
	\begin{align*}
		I_1 &:= I \cap \{ v \in \mathbb{N} : 1 \leq v \leq m \} \\
		I_2 &:= I \cap \{ v \in \mathbb{N} : m+1 \leq v \leq d+m \} \\
		I_3 &:= I \cap \{ v \in \mathbb{N} : d+m+1 \leq v \leq d+2m \}
	\end{align*}
	Then
	\begin{align*}
		(*) &= \sum_{v_1 \in I_1} k_{v_1-i-m} p_{m+1-v_1} +
		\sum_{v_2 \in I_2} k_{v_2-i-m} p_{v_2-m} +
		\sum_{v_3 \in I_3} k_{v_3-i-m} p_{2d+m+1-v_3}
	\end{align*}
	We apply a variable transform on each sum.
	\begin{align*}
		w_1 = w_1(v_1) := m-v_1+1 &\implies v_1=m+1-w_1 \\
		&\implies \widetilde{I_1} := w_1(I_1) \\
		& \phantomrel{\implies \widetilde{I_1} :}
		= \{ w \in \mathbb{N} : 1-m-i \leq w \leq 1+m-i \} \\
		& \phantomrel{\implies \widetilde{I_1} :=}
		\cap \{ w \in \mathbb{N} : 1 \leq w \leq m \}  \\
		& \phantomrel{\implies \widetilde{I_1} :}
		= \{ w \in \mathbb{N} : 1 \leq w \leq 1+m-i \} \\[20pt]
		%
		w_2 = w_2(v_2) := v_2-m &\implies v_2=w+m \\
		&\implies \widetilde{I_2} := w_2(I_2) \\
		&\phantomrel{\implies \widetilde{I_2} :}
		= \{ w \in \mathbb{N} : i-m \leq w \leq i+m \} \\
		& \phantomrel{\implies \widetilde{I_2} :=}
		\cap \{ w \in \mathbb{N} : 1 \leq w \leq d \}  \\
		&\phantomrel{\implies \widetilde{I_2} :}
		= \{ w \in \mathbb{N} : i-m \leq w \leq i+m \} \\[20pt]
		%
		w_3 = w_3(v_3) := 2d+m+1-v_3 &\implies v_3=2d+m+1-w \\
		&\implies \widetilde{I_3} := w_3(I_3) \\
		& \phantomrel{\implies \widetilde{I_3} :}
		= \{ w \in \mathbb{N} : 2d+1-m-i \leq w \leq 2d+m+1-i \} \\
		& \phantomrel{\implies \widetilde{I_3} :=}
		\cap \{ w \in \mathbb{N} : d+1-m \leq w \leq d \} \\
		& \phantomrel{\implies \widetilde{I_3} :}
		= \{ w \in \mathbb{N} : 2d+1-m-i \leq w \leq d \}
	\end{align*}
	Then
	\begin{align*}
		(*) = \sum_{w_1 \in \widetilde{I_1}} k_{1-i-w_1} x_{w_1} +
		\sum_{w_2 \in \widetilde{I_2}} k_{w_2-i} x_{w_2} +
		\sum_{w_3 \in \widetilde{I_3}} k_{2d+1-i-w_3} x_{w_3}
	\end{align*}
	Let $ 1 \leq j \leq d$. Then for the Jacobian holds
	\begin{align*}
		\hspace{6em}&\hspace{-6em}
		\lb conv1d \lsb k \rsb \circ symmetric\_pad \lsb m \rsb \rb_{ij}
		= \deldel{p_j} \lb conv1d \lsb k \rsb \circ symmetric\_pad \lsb m \rsb \rb_i \\
		&= \underbrace{\lb \begin{cases}
			k_{1-i-j} & : j \in \widetilde{I_1} \\
			0 & : else
		\end{cases} \rb}_{:= K_1(i,j)} + \underbrace{\lb \begin{cases}
			k_{j-i} & : j \in \widetilde{I_2} \\
			0 & : else
		\end{cases} \rb}_{:= K_2(i,j)} + \underbrace{\lb \begin{cases}
			k_{2d+1-i-j} & : j \in \widetilde{I_3} \\
			0 & : else
		\end{cases} \rb}_{:= K_3(i,j)} \\
		&= K_1(i,j) + K_2(i,j) + K_3(i,j)
	\end{align*}
	We have
	\begin{align*}
		 K_1(i,j) &= \begin{cases}
			k_{1-i-j} & : j \in \widetilde{I_1} \\
			0 & : else
		\end{cases} = \begin{cases}
			k_{1-i-j} & : 1 \leq j \leq 1+m-i \\
			0 & : else
		\end{cases} \\[7pt]
		&= \begin{cases}
			k_{1-i-j} & : j \leq 1+m-i \\
			0 & : else
		\end{cases} = \begin{cases}
			k_{1-i-j} & : j+i \leq 1+m \\
			0 & : else
		\end{cases} \\[25pt]
		%
		K_2(i,j) &= \begin{cases}
			k_{j-i} & : j \in \widetilde{I_2} \\
			0 & : else
		\end{cases} = \begin{cases}
			k_{j-i} & : i-m \leq j \leq i+m \\
			0 & : else
		\end{cases} = \begin{cases}
			k_{j-i} & : \abs{j-i} \leq m \\
			0 & : else
		\end{cases} \\[25pt]
		%
		K_3(i,j) &= \begin{cases}
			k_{2d+1-i-j} & : j \in \widetilde{I_3} \\
			0 & : else
		\end{cases} = \begin{cases}
			k_{2d+1-i-j} & : 2d+1-m-i \leq j \leq d \\
			0 & : else
		\end{cases} \\[7pt]
		&= \begin{cases}
			k_{2d+1-i-j} & : 2d+1-m-i \leq j \\
			0 & : else
		\end{cases}
		= \begin{cases}
			k_{2d+1-i-j} & : 2d+1-m \leq j+i \\
			0 & : else
		\end{cases}
	\end{align*}
	We show that $K_l$, $l=1,2,3$ are symmetric, i.e. $K_l(j,i) = K_l(i,j)$ for
	$l=1,2,3$. This implies that the Jacobian is symmetric.
	\begin{equation*}
		K_1(j,i) = \begin{cases}
			k_{1-j-i} & : i+j \leq 1+m \\
			0 & : else
		\end{cases} = \begin{cases}
			k_{1-i-j} & : j+i \leq 1+m \\
			0 & : else
		\end{cases} = K_1(i,j) \\[20pt]
	\end{equation*}
	\begin{equation*}
		K_2(j,i) = \begin{cases}
			k_{i-j} & : \abs{i-j} \leq m \\
			0 & : else
		\end{cases} \stackrel{k \text{ symmetric}}{=} \begin{cases}
			k_{j-i} & : \abs{j-i} \leq m \\
			0 & : else
		\end{cases} = K_2(i,j) \\[20pt]
	\end{equation*}
	\begin{equation*}
		K_3(j,i) = \begin{cases}
			k_{2d+1-j-i} & : 2d+1-m \leq i+j \\
			0 & : else
		\end{cases} = \begin{cases}
			k_{2d+1-i-j} & : 2d+1-m \leq j+i \\
			0 & : else
		\end{cases} = K_3(i,j)
	\end{equation*}
	Thus symplecticity follows with \cref{jacobi_symmetric}.
\end{proof}

\subsection{Activation layers}

$V(p) = 1_n^T(\int \sigma)(p)$

\begin{equation*}
	\grad{V}(p) = \lb\jac{V}{p}(p)\rb^T = \lb1_n^Tdiag\lb\sigma(p)\rb\rb^T
	= diag(\activation(p)) 1_n
	= \activation(p)
\end{equation*}

\todo{Differentiate between $a \in \mathbb{R}$ and $a \in \mathbb{R}^d$?}

\subsection{Gradient layers}

\begin{equation*}
	\mathcal{G}_{up} \qpvec := \uppersympop{\hat{\activation}} := \begin{pmatrix}
		q + K^T diag(a) \activation(Kp + b) \\
		p
	\end{pmatrix}
\end{equation*}
\begin{equation*}
	\mathcal{G}_{low} \qpvec := \lowersympop{\hat{\activation}} := \begin{pmatrix}
		q \\
		K^T diag(a) \activation(Kq + b) + p
	\end{pmatrix}
\end{equation*}
where $K \in \mathbb{R}^{n \times d}$ and $a,b \in \mathbb{R}^n$
are learnable parameters.
$n \in \mathbb{N}$ denotes the width of the gradient layer and can be chosen freely.

\newtheorem*{glayer}{Lemma}
\begin{glayer}
	Given an activation function $\sigma \in C^1$ the gradient layers $\mathcal{G}_{\{up,low\}}$ 
	are symplectic.
\end{glayer}
\begin{proof}
	Let $\mathcal{A}: \mathbb{R} \to \mathbb{R}$ be the antiderivative of $\activation$, 
	i.e. $\dd{x} \mathcal{A} = \activation$ and
	$V(p) := \onevec{n}^Tdiag(a)\mathcal{A}(Kp+b)$.

	Then $V \in C^2(\mathbb{R}^d, \mathbb{R})$ and
	\begin{align*}
		\nabla V(p) &= \left(\jac{V}{p}(p)\right)^T \\
		&= \left(1_n^Tdiag(a)diag\left(\sigma(Kp+b)\right)K\right)^T \\
		&= K^Tdiag(a)diag\left(\sigma(Kp+b)\right)1_n \\
		&= K^Tdiag(a)\sigma(Kp+b)
	\end{align*}

	With \cref{gradient_corollary} follows that the Gradient layers $\mathcal{G}_{\{up,low\}}$
	are symplectic.
\end{proof}

\todo{For high dimensions $K$ and $K^T$ can be implemented as transposed convolution
and convolution operators.}

\subsection{Normalization}

\todo{Introduce batch normalization from literature.}

\todo{Normalized activation layers}

\todo{Normalized gradient layers}

\section{Relation to geometric integrators}

\section{Numerical experiments}

\subsection{Low-dimensional systems}

\subsubsection{Harmonic Oscillator}

\subsubsection{Simple Pendulum}

\subsection{High-dimensional systems}

\todo{Describe Hamiltonian PDEs. A PDE describes system locally, thus it makes sense to share parameters, for
example in the form of CNNs etc.}

Assumptions:
\begin{itemize}
	\item Original PDE has constant coefficients \todo{Or rather should not depend on $x$ itself?}
	\item Equidistant grid points
\end{itemize}

\todo{Explain why we need assumptions.}

\subsubsection{Linear wave equation}

\todo{The choice of basis for the symmetric convolution kernels is critical (at least for fast convergence)!}

\subsubsection{Sine-Gordon}

\todo{Bias important if non-zero Dirichlet boundaries, but only at borders! // Padding modes}

\todo{Activation layer with sin works well if $a \in \mathbb{R}$, i.e. activation function is applied
isotropic.}

\todo{Maybe compare different optimizers (SGD, Adam, AdamW, amsgrad, ...)}

\section{Implementation}

\todo{Describe use of PyTorch, TensorBoard, Docker, MLFlow, Unit Tests, ...}

%----------------------------------------------------------------------------------------
%
% Resume
%
%----------------------------------------------------------------------------------------
\clearpage\newpage\null %empty page
\newpage
\section{R\'esum\'e}
\subsection{Summary and conclusion}

\subsection{Outlook}

%----------------------------------------------------------------------------------------
%
% Appendices
%
%----------------------------------------------------------------------------------------
\clearpage\newpage\null %empty page
\newpage
\begin{appendices}
\addtocontents{toc}{\protect\setcounter{tocdepth}{1}}
\section{First Appendix Section}

\newpage~\newpage
\section{Declaration of authorship}

\vspace{3cm}

\begin{table}[h!]
\centering
\begin{tabular}{|p{13cm}|}
\hline\\
	I hereby certify
	\begin{enumerate}
		\item that this thesis has been composed by me and is based on my own work, unless stated otherwise,
		\item that all direct or indirect sources used are acknowledged as references and all extracts from work of others, either verbatim or in spirit, are stated as such,
		\item that neither the thesis itself nor parts of this thesis have been part of another examination procedure,
		\item that neither the thesis itself nor parts of this thesis have been published and
		\item that all copies of this thesis, either digital or printed, coincide.
	\end{enumerate}
	Therewith, this declaration of authorship is in accordance with the examination regulations from 29th July 2013 of the master's program \emph{Simulation Technology} of the University of Stuttgart.\\\\
\hline
\end{tabular}
\end{table}

\vspace{4cm}
\hrulefill\\
Name
\hspace{7cm}
Date, City, Signature
\end{appendices}
%----------------------------------------------------------------------------------------
%
% BIBLIOGRAPHY
%
%----------------------------------------------------------------------------------------
\clearpage\newpage\null %empty page
\newpage
\addtocontents{toc}{\protect\setcounter{tocdepth}{1}}
\addcontentsline{toc}{section}{References}
\bibliographystyle{abbrvnat}
\bibliography{../../literature/references.bib}

\end{document}
